\section{区际司法协助}
\begin{enumerate}
  \item 重点考查对象，尤其是涉港澳区际司法协助。香港暴乱。
  \item 文书送达、调取证据、判决的跨地区承认和执行。两两分开。
  \item 文书送达。港澳 7 种，涉台 8 种，域外 9 种。
        \begin{enumerate}
          \item 涉港澳台\zd{委托送达}途径有双边安排\zd{不叫双边条约}。期限 2 月，收到起算。\[
  \begin{tabular}{cccc}
    \toprule
                        & 涉港            & 涉澳                            & 涉台                               \\
    \midrule
    \multirow{2}{*}{机构} & 高院 \zyjt 高等法院 & 高院、最高法\bs{授权的}中基 \zyjt   终审法院 & \multirow{2}{*}{高院 \zyjt   有关法院} \\ \cmidrule(lr){2-3}
                        & 最高法 \yjt 高等法院 & 最高法 \zyjt 终审法院                &                                  \\
    \bottomrule
  \end{tabular}
\]

          \item 涉台送达可指定代收人。
          \item 认定送达。
                \begin{enumerate}
                  \item 签收行为（送达回证 \hz 邮件回执）。
                  \item 受送达人某些行为。提到文书内容、按文书要求履行义务等。
                  \item 留置（\zd{对象明确且在内地}）。
                  \item 公告期满 3 月。
                \end{enumerate}
          \item 邮件期满 3 月音讯全无视为未送达。
        \end{enumerate}
  \item 调取证据。没有和台湾达成取证的安排。
        \begin{enumerate}
          \item 我方主体不变，对方仅香港变，\zd{香港政务司行政署}，委托书加盖\zd{香港高等法院签章}。
          \item 期限。涉港 6 个月，涉澳 3 个月。
          \item 只能是司法证据。
          \item \zd{委托方可以请求直接参与取证，受托方可以拒绝/同意 \yjt 应当通知取证的时间地点}。\al{直接向证人提问 \dui}。
                \begin{kk}{新的成果}
                  \begin{enumerate}
                    \item 最高法授权的中基直接与澳终审法院对接文书和调取证据。
                    \item 增加音视频取证。委托方请求，证人、鉴定人同意，受托方\zd{可以}协助安排。
                    \item 委托书上法院印章和法官签名具有同等效力。
                  \end{enumerate}
                \end{kk}
        \end{enumerate}
  \item 判决、仲裁\zd{认可}与执行。
        \begin{enumerate}
          \item 内地和港澳双边（即任意法院、仲裁机构），和台湾单边。\al{国务院港澳办提供的名单内的仲裁机构}。
          \item 找有管辖权的\zd{中院}。
          \item 香港\zd{婚姻家事找区域，其他判决找高等}。判仲执行均允许拆分，不超过总额.
          \item 澳门\zd{认可、申请执行向中院提 \yjt 初级法院执行}，仅仲裁允许拆。
          \item 期限、程序、费用和方式依据\zd{被请求}方法律规定。
          \item 对内裁定不服向上一级法院复议；对港澳不服可上诉。\zd{不予受理、驳回起诉、管辖权异议}裁定可上诉。
          \item 同时申请执行。\zd{同时两地就澳判不能，但同时内地都不能，谁先立案谁受理}。
          \item 台湾。6 个月审结。
        \end{enumerate}

        \begin{kk}{港澳仲裁裁决的相互协助保全}
          \begin{enumerate}
            \item 范围。证据、行为和财产保全。
            \item 内地出中院（但只能找一个申请），香港高等法院，澳门初级法院。
            \item 保全的程序、费用、救济和方式依据\zd{被请求}方法律规定。
          \end{enumerate}
        \end{kk}
  \item 司法协助（涉外和区际）共同规则。
        \begin{enumerate}
          \item \zd{专属管辖可以成为拒绝承认/认可法院判决的理由}，但不能成为拒绝其他司法协助的理由。
          \item 应支付执行费用；文书送达、证据无须支付特别费用，但须承担法院实际产生的特别费用。
          \item 向中国内地法院请求司法协助原则上都应当提交有关文件的中文译本。
          \item 外国法院依据国际条约送达文书 \hz 调取证据时，可附条约规定的第三方文本。
        \end{enumerate}
\end{enumerate}
